% appendix-a.tex - Exit Codes and Status Reference

\chapter{Exit and Status Code Reference}
\label{appendix:codes}

This appendix provides a complete reference of all exit codes and status values used by \SC{}.

\section{Command Exit Codes}

The following exit codes are returned by \SC{} commands in the response data field when a command is rejected:

\begin{longtable}{clp{7cm}}
    \toprule
    \textbf{Code} & \textbf{Constant} & \textbf{Description} \\
    \midrule
    \endfirsthead
    \toprule
    \textbf{Code} & \textbf{Constant} & \textbf{Description} \\
    \midrule
    \endhead
    \midrule
    \multicolumn{3}{r}{\textit{Continued on next page}} \\
    \endfoot
    \bottomrule
    \endlastfoot
    \hex{00} & SUCCESS & Process accepted without error \\
    \hex{01} & INVALID\_ARGS & Invalid arguments to command; unrecognized DR name, unknown copy ID, or unknown MIB entry \\
    \hex{02} & OTHER\_ERROR & Other error running command; general error condition such as failure to queue a copy or delete operation \\
    \hex{03} & BLOCKING\_OP & Blocking operation in progress; INI or SHT command is currently executing and prevents new INI/SHT operations \\
    \hex{04} & NOT\_INITIALIZED & Subsystem needs to be initialized; system is in SHUTDWN state and requires INI before accepting data management commands \\
\end{longtable}

\section{System Status Values}

The 7-character status strings reported via \mibentry{SUMMARY}:

\begin{longtable}{lp{10cm}}
    \toprule
    \textbf{Status} & \textbf{Description} \\
    \midrule
    \endfirsthead
    \toprule
    \textbf{Status} & \textbf{Description} \\
    \midrule
    \endhead
    \midrule
    \multicolumn{2}{r}{\textit{Continued on next page}} \\
    \endfoot
    \bottomrule
    \endlastfoot
    \status{BOOTING} & System is initializing. Occurs after INI command is received and while initialization sequence is in progress. Data management commands are not accepted. \\
    \status{NORMAL } & System operating normally. All DR queue managers running, station monitor active, and system ready to accept data management commands. \\
    \status{SHUTDWN} & System has been shut down. Occurs after SHT command completes. All monitoring and queue processing threads stopped. Pending queue items are preserved in the database. Requires INI to restart. \\
\end{longtable}

\section{Response Code Characters}

Single-character response codes in packet byte 38 are listed in Table~\ref{tab:response-chars}.

\begin{table}[htbp]
    \centering
    \caption{Response Code Characters}
    \label{tab:response-chars}
    \begin{tabular}{clp{8cm}}
        \toprule
        \textbf{Code} & \textbf{Meaning} & \textbf{Description} \\
        \midrule
        \code{A} & Accepted & Command was accepted and executed successfully \\
        \code{R} & Rejected & Command was rejected; see response data for exit code and reason \\
        \bottomrule
    \end{tabular}
\end{table}

\section{Queue Item Statuses}

Status values stored in the \code{queue\_items.status} column of the SQLite database:

\begin{longtable}{lp{10cm}}
    \toprule
    \textbf{Status} & \textbf{Description} \\
    \midrule
    \endfirsthead
    \toprule
    \textbf{Status} & \textbf{Description} \\
    \midrule
    \endhead
    \midrule
    \multicolumn{2}{r}{\textit{Continued on next page}} \\
    \endfoot
    \bottomrule
    \endlastfoot
    \code{pending} & Item is in the queue waiting to be processed. This is the initial status for all new items and the status items return to after a failed attempt with retries remaining. \\
    \code{processing} & Item has been dequeued and is currently being copied or deleted. If found on startup (indicating an interrupted operation), the item is automatically reset to \code{pending}. \\
    \code{completed} & Transfer completed successfully. The source file is eligible for deletion during the next purge cycle. These records accumulate until the purge threshold is reached. \\
    \code{failed} & Transfer has permanently failed due to exhausted retries or missing source file. A failure notification will be sent via email and the record will be purged. \\
\end{longtable}

\section{Copy Operation Statuses}

Status strings returned by the \code{InterruptibleCopy} engine and visible in \mibentry{ACTIVE\_STATUS} MIB entries:

\begin{longtable}{lp{9cm}}
    \toprule
    \textbf{Status} & \textbf{Description} \\
    \midrule
    \endfirsthead
    \toprule
    \textbf{Status} & \textbf{Description} \\
    \midrule
    \endhead
    \midrule
    \multicolumn{2}{r}{\textit{Continued on next page}} \\
    \endfoot
    \bottomrule
    \endlastfoot
    \code{active} & Copy is currently running \\
    \code{active/started for ...} & Copy has just been started (includes source and destination) \\
    \code{active/resumed for ...} & Copy has been resumed after a pause (includes source and destination) \\
    \code{paused} & Copy has been paused (rsync process killed, ready to resume) \\
    \code{paused for ...} & Copy is paused (includes source and destination) \\
    \code{queued for ...} & Copy is waiting in the queue (includes source and destination) \\
    \code{complete} & Copy completed successfully \\
    \code{canceled} & Copy was canceled via SCN command \\
    \code{error: ...} & Copy failed (includes error details from rsync stderr) \\
    \code{error: too soon to retry} & Retry interval has not elapsed since last attempt \\
\end{longtable}

\section{Error Response Format}

Error responses are returned with an \code{R} (Rejected) response code in byte 38 of the response packet. The response data contains:

\begin{enumerate}
    \item Exit code in hexadecimal (e.g., \hex{04})
    \item Exclamation point delimiter
    \item Command name and colon separator
    \item Descriptive error message from the \mibentry{LASTLOG} state
\end{enumerate}

Example: \code{0x01! SCP: Invalid arguments to command - unknown data recorder DR7}
